% ====================================================================
\section{Discussion}\label{sec:discussion}
% ====================================================================
\begin{table}[t]
  \centering\small
  \caption{Results at a glance ($1\sigma$; one-sided systematics in parentheses).}
  \label{tab:results}
  \begin{tabularx}{\textwidth}{@{}l l >{\raggedright\arraybackslash}X@{}}
    \toprule
    Quantity & Result & Comparison \\
    \midrule
    Hall probe at sample position & \qty{-8.5}{\percent} (analog), \qty{-2.3}{\percent} (display) & DPPH as standard \\
    Hysteresis of the yoke & \qtyrange{14}{24}{mT}, residual field \qtyrange{17}{24}{mT} & -- \\
    DPPH: $g$ with display & \num{2.052} & \num{2.0036}~\cite{poole1996} \\
    DPPH: $\Bpp$, $T_2$ & \qty{0.240(15)}{mT}, \qty{27(2)}{ns} (up to \qty{35}{ns}) & \qtyrange{0.15}{0.47}{mT}~\cite{poole1996} \\
    Modulation at \qty{1}{V} & $\Bm<\qty{0.059}{mT}$ & recommended: $\Bpp/5$ to $\Bpp/3$ \\
    Power dependence & $S_\mathrm{pp}\propto P^{\num{0.98(3)}}$ & square-law diode: $P^1$ \\
    Bridge leak & $\rho\approx\num{0.26}$, $\cos\psi\approx\num{-0.78}$ & -- \\
    Dispersion, model-free & $D=\qty{-84.3(7)}{\degree}$ & model fit: \ang{-87.9} \\
    Guide wavelength & $\lamg=46^{+8}_{-6}\,\unit{mm}$ & WR-90: \qty{44.9}{mm} \\
    Cutoff frequency & $\nu_\mathrm c=6.7^{+0.8}_{-1.1}\,\unit{GHz}$ & WR-90: \qty{6.56}{GHz} \\
    \ce{MnCl2}: $\Bhalf$, $T_2$ & \qty{88.8(23)}{mT}, \qty{62(2)}{ps} & without exchange: about \qty{250}{mT} \\
    \ce{MnCl2}: $g$ & \num{2.058(11)} (\num{-0.025}) & $\num{2.001}+\Delta g$, $\Delta g=\numrange{0.028}{0.055}$ \\
    \bottomrule
  \end{tabularx}
\end{table}

\paragraph{Magnetic field.}
The magnetic-field calibration was the largest source of uncertainty in this experiment (\cref{tab:results} collects
all results). At the sample position the Hall probe reads too low, and only the DPPH calibration turns the control voltage into a
reliable field. Next time I would first move the probe to the height and orientation with the highest reading, as the
course supervisor recommends. Any deficit left after that would point to the gaussmeter calibration. The DPPH
calibration holds only for scans run exactly like the DPPH scans. Because the field lags behind the set point, starting
just \qty{75}{mV} earlier shifts the line by \qty{1.2}{mT}, almost five times the DPPH line width. This
causes the three large systematics of the report: the DPPH width may be overestimated by up to \qty{22}{\percent}, the
field reference for \ce{MnCl2} is uncertain, and the \ce{MnCl2} scan may have run on a different hysteresis branch. A
small marker sample would avoid all three, for example a grain of DPPH next to the sample in the cavity. It would appear
in every scan, with the same magnetic history as the measurement. Without a marker, every scan should start from
\qty{0}{V} with the same dwell time and sweep rate, and every deviation should be noted in the lab notebook.

\paragraph{Modulation and DPPH.}
For the modulation amplitude I only have an upper bound. A small pick-up coil at the sample position would have measured
it directly. According to the course supervisor (private communication), such a coil is available in the lab, in the
drawer with the waveguide parts; I did not ask for it during the measurement: with 100 turns on \qty{10}{mm^2} and $\Bm=\qty{0.05}{mT}$ it would give about \qty{40}{\micro V}, easy to
measure with the lock-in. The value of $T_2$ rests on the convolution model, because the large powder sample sees a
field distribution of \qty{0.82}{mT}, almost four times the HWHM of the line. With only a few grains in the center of the
cavity, a plain Lorentzian fit would give $T_2$ directly.

\paragraph{Bridge.}
The lab manual uses the main attenuator for the power dependence and the attenuator in the reference arm for
balancing the bridge, and my measurements support this choice. With X375A the signal follows the power linearly, as
expected for a diode at low power, while X382A changes only the reference. The Hilbert-transform fit shows without any
line model that absorption and dispersion obey the Kramers--Kronig relations within the noise. The leak of
the cavity makes phases comparable only within one measurement block and limits the accuracy of $\lamg$ to about
\qty{15}{\percent}. Several plunger positions in quick succession would give a more accurate value. The complex amplitude
then moves on a circle around the fixed leak contribution, and the rotation angle per millimeter can be fitted
independently of it.

\paragraph{\texorpdfstring{\ce{MnCl2}}{MnCl2} and copper.}
The \ce{MnCl2} line shows the typical features of exchange narrowing. It is Lorentzian, has no hyperfine structure, and is
only a third as wide as the dipolar fields would make it without exchange. If the deviation of its $g$-factor from 2 comes
from the large width, as I suspect, it should almost vanish at higher frequency. The copper measurement has to be repeated
with a properly tuned source, reading the frequency from the counter before and after each scan.
